% Generated draft: requires mathematical and source review.
\documentclass[11pt]{article}
\usepackage[margin=25mm]{geometry}
\usepackage{fontspec}
\setmainfont{Latin Modern Roman}
\usepackage{amsmath,amssymb,mathtools,unicode-math}
\setmathfont{Latin Modern Math}
\usepackage{enumitem,needspace,xurl,hyperref}
\setlength{\parindent}{0pt}
\setlength{\parskip}{5pt}
\setlength{\emergencystretch}{3em}
\newcommand{\N}{\mathbb{N}}
\newcommand{\Z}{\mathbb{Z}}
\newcommand{\Q}{\mathbb{Q}}
\newcommand{\R}{\mathbb{R}}
\newcommand{\C}{\mathbb{C}}
\newcommand{\F}{\mathbb{F}}
\newcommand{\A}{\mathbb{A}}
\newcommand{\PP}{\mathbb{P}}
\newcommand{\E}{\mathbb{E}}
\newcommand{\eps}{\varepsilon}
\newcommand{\dd}{\,\mathrm{d}}
\begin{document}
\section*{UnsolvedMath ID: 30001610}
\textbf{Source problem code:} catalogue entry
\subsection{Definitions and mathematical statement}
The supplied problem asks for a characterization of positive weights \(\omega:\Z\to(0,\infty)\) having bounded successive ratios for which
\[
\exists x\in\ell^p_\omega(\Z)\qquad
\overline{\operatorname{span}\{S^n x:n\in\Z_{\geq0}\}}^{\,\ell^p_\omega(\Z)}
=\ell^p_\omega(\Z).
\]
Here completeness is interpreted as density of the linear span of the right translates. With the usual coordinate convention, right translation is \((Sx)(k)=x(k-1)\) for every \(k\in\Z\). Thus the target is a necessary and sufficient condition on \(\omega\), for each permitted \(p\), for the existence of such a cyclic vector \(x\).

The source record does not define the weighted norm, scalar field, or precise ratio hypothesis. Consequently, the displayed target requires the following conventions to be resolved before it becomes a uniquely specified mathematical problem. For a chosen scalar field \(\mathbb K\in\{\R,\C\}\), a real exponent \(1\leq p<\infty\), and a positive sequence \(v\), define
\[
\ell^p(v;\Z;\mathbb K)
=\left\{x:\Z\to\mathbb K:
\sum_{k\in\Z}|x(k)|^p v(k)<\infty\right\},
\qquad
\|x\|_{p,v}=\left(\sum_{k\in\Z}|x(k)|^p v(k)\right)^{1/p}.
\]
Two possible interpretations of the source notation are \(\ell^p_\omega=\ell^p(\omega;\Z;\mathbb K)\) and \(\ell^p_\omega=\ell^p(\omega^p;\Z;\mathbb K)\); the supplied record does not select one. Likewise, bounded successive ratios could mean
\[
\sup_{k\in\Z}\frac{\omega(k+1)}{\omega(k)}<\infty,
\]
or the stronger two-sided condition
\[
\exists C\geq1\ \forall k\in\Z:\qquad
C^{-1}\leq\frac{\omega(k+1)}{\omega(k)}\leq C.
\]
No choice between these hypotheses is asserted here. The source writes \(p\geq1\) and does not explicitly address an \(\ell^\infty\) endpoint. The finite-exponent definitions above therefore specify conditional formulations only; the intended conventions and any endpoint formulation require source confirmation.
\subsection{Short English statement}
Characterize the positive weights on \(\Z\), subject to the intended bounded-ratio hypothesis, for which some vector has right translates whose linear span is dense in the weighted \(\ell^p\) space. The weighted norm, scalar field, exact ratio condition, and any \(p=\infty\) endpoint remain unspecified in the supplied record.
\subsection{Sources}
\begin{itemize}
\item Supplied local UnsolvedMath registry record, permanent ID \(30001610\), exact source problem code \texttt{catalogue entry}, title: \emph{Cyclic Right Translates in Weighted \(\ell^p\) Spaces}. The record lists proposed year 2010 and supplies no named proposer.
\item Supplied report citation: \emph{Operator Theory and Harmonic Analysis (2011)}, Oberwolfach Reports, pp.~2813--2875, DOI \texttt{10.4171/owr/2010/49}. \url{https://doi.org/10.4171/owr/2010/49}. The background describes Oberwolfach Reports as documenting mathematical workshops and their open-problem discussions.
\item Supplied literature source labeled ``Abakumov Atzmon and Grivaux problem statement'': \url{https://ems.press/content/serial-article-files/46308}. This label is retained as supplied attribution; it does not establish individual proposership.
\item Supplied literature source labeled ``original problem report'': \url{https://doi.org/10.4171/owr/2010/49}.
\item Literature assessment dated exactly \texttt{2026-08-21}: the supplied curation labels the problem open, reports that the report explicitly describes the characterization as still open, and reports partial results but no general characterization found in its search. This is supplied dated curation, not an independently verified assessment. The background expressly calls for independent verification of time-sensitive status and cited claims.
\end{itemize}
\end{document}
